\chapter{Craniotomy}
\section*{Overview}
Craniotomy temporarily removes a section of skull to access the brain for diagnosis or treatment.
\section*{Why This Test or Procedure Is Ordered}
It may be required for brain tumors, hemorrhage, vascular malformations, aneurysms, trauma, epilepsy surgery, or raised intracranial pressure.
\section*{How This Test or Procedure Is Performed}
Under general anesthesia, a bone flap is created and replaced after the intracranial procedure. Image guidance, neurophysiologic monitoring, or awake techniques may be used for selected cases.
\section*{Risks and Recovery}
Bleeding, infection, seizures, stroke, brain swelling, neurologic deficits, cerebrospinal-fluid leak, and anesthesia complications are possible. Recovery depends on the underlying condition and neurologic findings.
\section*{References}
\begin{enumerate}\item MedlinePlus. Craniotomy. U.S. National Library of Medicine; 2025.\item Current Medical Diagnosis and Treatment. McGraw-Hill; 2025.\end{enumerate}
\clearpage
